\documentclass[12pt, letterpaper]{article}
\usepackage[utf8]{inputenc}
\usepackage[spanish]{babel}
\usepackage{amsmath, amssymb, amsthm, bm}
\usepackage{graphicx}
\usepackage{array}
\usepackage{enumitem}
\usepackage{float}
\usepackage{listings}
\usepackage{url}
\usepackage{xcolor}
\usepackage{tabularx}
\usepackage{booktabs}
%% Sets page size and margins
\usepackage[top=2.5cm,bottom=2.5cm,left=2cm,right=2cm]{geometry}
\usepackage{setspace}
\onehalfspacing
\usepackage{gensymb}
\usepackage{fancyhdr}
\pagestyle{fancy}
\fancyhead[L]{Compiladores}
\fancyhead[C]{Práctica 6}
\fancyhead[R]{Facultad de Ciencias UNAM}
\fancyfoot[C]{\thepage}
\renewcommand{\headrulewidth}{0.4pt}
\renewcommand{\footrulewidth}{0.4pt}

\lstdefinestyle{bashstyle}{
    backgroundcolor=\color{gray!10},
    basicstyle=\small\ttfamily,
    breaklines=true,
    frame=single,
    frameround=tttt
}

\lstset{
    inputencoding=utf8,
    extendedchars=true,
    literate=
        {á}{{\'a}}1
        {é}{{\'e}}1
        {í}{{\'i}}1
        {ó}{{\'o}}1
        {ú}{{\'u}}1
        {Á}{{\'A}}1
        {É}{{\'E}}1
        {Í}{{\'I}}1
        {Ó}{{\'O}}1
        {Ú}{{\'U}}1
        {ñ}{{\~n}}1
        {Ñ}{{\~N}}1
        {ü}{{\"u}}1
        {Ü}{{\"U}}1
}


\title{
    \vspace*{-2.5cm}
    \setlength{\parindent}{0cm} %Indentation 0
    \rule{\linewidth}{0.1mm}
    \vspace*{-1.4cm}
    \begin{center}
        \begin{minipage}{2.5cm}
            \begin{center}
            \includegraphics[height=3.4cm]{images/Logo_UNAM.png}
            \end{center}
        \end{minipage}\hfill
        %--------------
        \hspace{0.3cm}
        \begin{minipage}{11cm}
            \begin{center}
            \normalfont \small \textsc{UNIVERSIDAD NACIONAL AUTÓNOMA DE MÉXICO \\ FACULTAD DE CIENCIAS \\ Compiladores } \\
    		\normalsize \textbf{Práctica 6}\\
            \normalsize \textsc{Parser Canónico LR(1) y Compactación LALR(1)}\\
            \end{center}
        \end{minipage}\hfill
        %--------------
        \begin{minipage}{3cm}
            \begin{center}
            \includegraphics[height=3.6cm]{images/Logo_FC.png}
            \end{center}
        \end{minipage}
    \end{center}
    \vspace*{-0.25cm}
    \rule{\linewidth}{0.1mm}
    \vspace*{0.5cm}
    {\normalsize \textbf{Integrantes:} \par}
    \vspace{-0.5cm}
    \begin{tabular}{{l@{\hspace{1cm}}r}}
        \vspace*{-0.5cm}
        \small Escobar Gonzalez Isaac Giovani & \small 321336400 \\
        \vspace*{-0.5cm}
        \small Garduño Escobar Kevin Jonathan & \small 321070629 \\
        \vspace*{-0.5cm}
        \small González Gutiérrez Antonio Tonatiuh & \small 321195476 \\
        \vspace*{-0.5cm}
        \small Zaldivar Alanis Rodrigo & \small 424029605
    \end{tabular}
    \date{}
}

\begin{document}
\maketitle
\vspace*{-2.5cm}

\subsection*{Diseño de Solución}

\newpage

\subsection*{Tabla de Pruebas y Resultados}


\textbf{Comandos de compilación y ejecución:}

Construir la imagen del contenedor:
\begin{lstlisting}[style=bashstyle]
docker build -t parser-lr .
\end{lstlisting}
Ejecutar el motor de análisis:
\begin{lstlisting}[style=bashstyle]
docker run --rm parser-lr
\end{lstlisting}

\textbf{Enlace al repositorio público:} 
\url{https://github.com/KevinJGard/Compiladores}

\newpage
\subsection*{Validación de Trabajo Presencial / Reposición Teórica}

\begin{itemize}
    \item \textbf{Isaac Giovani Escobar González}:

    IMG

    \newpage

    \item \textbf{Kevin Jonathan Garduño Escobar}:
    
    IMG

    \newpage
    
    \item \textbf{Antonio Tonatiuh González Gutiérrez}:
    
    IMG

    \newpage
    
    \item \textbf{Rodrigo Zaldivar Alanis}:
    
    IMG
    
\end{itemize}

\newpage
\subsection*{Evaluación de la Práctica}
\noindent Grado de acuerdo con cada una de las siguientes afirmaciones utilizando una escala de 1 a 5, donde:
\begin{itemize}
    \item 1: Totalmente en desacuerdo
    \item 2: En desacuerdo
    \item 3: Neutral
    \item 4: De acuerdo
    \item 5: Totalmente de acuerdo
\end{itemize}

\begin{table}[H]
    \centering
    \begin{tabularx}{\textwidth}{|X|c|c|c|c|c|}
        \hline
        \textbf{Criterio / Afirmación} & \textbf{1} & \textbf{2} & \textbf{3} & \textbf{4} & \textbf{5} \\
        \hline
        El calcular manualmente el autómata me ayudó a entender la diferencia de potencia entre SLR(1) y LALR(1). &  &  & \checkmark &  &  \\
        \hline
        Programar el motor de la pila por separado de la tabla dejó más clara la arquitectura de un compilador. &  &  &  & \checkmark &  \\
        \hline
        El nivel de complejidad técnica fue adecuado y el no usar plantillas permitió mayor libertad creativa en el código. &  &  & \checkmark &  &  \\
        \hline
        Las pautas para analizar la gramática proporcionada fueron claras y guiaron bien el desarrollo. &  &  &  & \checkmark &  \\
        \hline
        Considero que separar el análisis teórico (a mano) del código fuente redujo mi frustración frente a la computadora. &  &  &  & \checkmark &  \\
        \hline
    \end{tabularx}
\end{table}

\subsection*{Comentarios y Retroalimentación}
\noindent Por favor, dedica un par de líneas para comentar:
\begin{itemize}
    \item ¿Qué aspecto de esta práctica fue tu favorito o te resultó más revelador?
    
    \textbf{Respuesta:}

    \item ¿Qué cambiarías o mejorarías de esta práctica para futuras generaciones?
    
    \textbf{Respuesta:}
    
\end{itemize}

\end{document}
